\chapter{NP-تمامیت}

استیون کوک و لئونید لوین در اوایل دهه ۱۹۷۰ به‌طور مستقل کشف کردند که بعضی مسئله‌های $\cNP$ پیچیدگی‌شان به پیچیدگی \textbf{کل} کلاس $\cNP$ وابسته است: اگر برای \textbf{یکی} از این مسئله‌ها الگوریتم چندجمله‌ای وجود داشته باشد، \textbf{همه} مسئله‌های $\cNP$ در زمان چندجمله‌ای حل می‌شوند. به این مسئله‌ها \term{NP-تمام}{NP-complete} می‌گوییم.

این پدیده از دو جهت مهم است:
\begin{itemize}
\item \textbf{نظری:} پژوهشگری که می‌خواهد ثابت کند $\cP=\cNP$، فقط کافی است برای یک مسئله NP-تمام الگوریتم چندجمله‌ای پیدا کند.
\item \textbf{عملی:} NP-تمام بودن یک مسئله شاهد محکمی است که الگوریتم چندجمله‌ای ندارد. پس به جای اتلاف وقت برای پیدا کردن چنین الگوریتمی، بهتر است سراغ راه‌حل‌های تقریبی یا حالت‌های خاص برویم.
\end{itemize}

\section{کاهش‌پذیری در زمان چندجمله‌ای}

\begin{definition}[تابع محاسبه‌پذیر در زمان چندجمله‌ای]
تابع $f:\Sigma^*\to\Sigma^*$ \textbf{محاسبه‌پذیر در زمان چندجمله‌ای} است اگر یک ماشین تورینگ با زمان چندجمله‌ای مثل $M$ وجود داشته باشد که روی هر ورودی $w$ با فقط $f(w)$ روی نوارش متوقف شود.
\end{definition}

\begin{definition}[کاهش چندجمله‌ای]
زبان $A$ \textbf{به‌طور چندجمله‌ای کاهش‌پذیر} به زبان $B$ است و می‌نویسیم $A\le_PB$، اگر تابعی محاسبه‌پذیر در زمان چندجمله‌ای مثل $f:\Sigma^*\to\Sigma^*$ وجود داشته باشد که برای هر $w$:
\[w\in A\iff f(w)\in B.\]
به $f$ \textbf{کاهش چندجمله‌ای} از $A$ به $B$ می‌گوییم.
\end{definition}

کاهش چندجمله‌ای همان کاهش نگاشتی فصل ۲۱ است، با این شرط اضافه که تبدیل باید \textbf{کارآمد} باشد. این کاهش آزمون عضویت در $A$ را به آزمون عضویت در $B$ تبدیل می‌کند.

\begin{theorem}\label{thm:polyred-p}
اگر $A\le_PB$ و $B\in\cP$، آنگاه $A\in\cP$.
\end{theorem}
\begin{proof}
فرض کنید $M$ الگوریتمی چندجمله‌ای برای $B$ و $f$ کاهش چندجمله‌ای از $A$ به $B$ باشد. الگوریتم $N$ برای $A$:
\begin{algo}{$N$}{روی ورودی $w$}
\begin{enumerate}
\item $f(w)$ را محاسبه کن.
\item $M$ را روی $f(w)$ اجرا کن و همان خروجی را بده.
\end{enumerate}
\end{algo}
اگر $w\in A$ آنگاه $f(w)\in B$ و $M$ می‌پذیرد؛ اگر $w\notin A$ آنگاه $f(w)\notin B$ و $M$ رد می‌کند. زمان: مرحله ۱ چندجمله‌ای است. چون $f$ در زمان چندجمله‌ای محاسبه می‌شود، طول $f(w)$ هم چندجمله‌ای بر حسب $\abs{w}$ است (ماشین در زمان $t$ حداکثر $t$ نماد می‌نویسد). $M$ روی ورودی‌ای با طول چندجمله‌ای در زمان چندجمله‌ای اجرا می‌شود و ترکیب دو چندجمله‌ای هم چندجمله‌ای است. پس $N$ در زمان چندجمله‌ای اجرا می‌شود.
\end{proof}

\section{NP-تمامیت}

\begin{definition}[NP-تمام]
زبان $B$ \textbf{NP-تمام} است اگر دو شرط را داشته باشد:
\begin{enumerate}
\item $B\in\cNP$، و
\item برای هر $A\in\cNP$: \ $A\le_PB$.
\end{enumerate}
\end{definition}

\begin{theorem}
اگر $B$ NP-تمام باشد و $B\in\cP$، آنگاه $\cP=\cNP$.
\end{theorem}
\begin{proof}
هر $A\in\cNP$ طبق شرط ۲ به‌طور چندجمله‌ای به $B$ کاهش می‌یابد و طبق قضیه \ref{thm:polyred-p}، $A\in\cP$. پس $\cNP\subseteq\cP$ و چون $\cP\subseteq\cNP$، برابرند.
\end{proof}

\begin{theorem}\label{thm:npc-chain}
اگر $B$ NP-تمام باشد و $B\le_PC$ برای یک $C\in\cNP$، آنگاه $C$ هم NP-تمام است.
\end{theorem}
\begin{proof}
$C\in\cNP$ را داریم؛ باید نشان دهیم هر $A\in\cNP$ به‌طور چندجمله‌ای به $C$ کاهش می‌یابد. چون $B$ NP-تمام است، $A\le_PB$ با کاهشی مثل $f$. همچنین $B\le_PC$ با کاهشی مثل $g$. کاهش‌های چندجمله‌ای قابل ترکیب‌اند: تابع $g\circ f$ (اول $f$، بعد $g$) در زمان چندجمله‌ای محاسبه می‌شود (چون طول $f(w)$ چندجمله‌ای است) و
\[w\in A\iff f(w)\in B\iff g(f(w))\in C.\]
پس $A\le_PC$.
\end{proof}

این قضیه راه اصلی اثبات NP-تمام بودن است: کافی است یک مسئله NP-تمام شناخته‌شده را به مسئله جدید کاهش دهیم. اما برای شروع به \emph{اولین} مسئله NP-تمام نیاز داریم.

\section{مسئله صدق‌پذیری}

اولین مسئله NP-تمام، مسئله \textbf{صدق‌پذیری} است. \textbf{متغیرهای بولی} مقدار $0$ یا $1$ می‌گیرند و با عملگرهای \lr{AND} ($\wedge$)، \lr{OR} ($\vee$) و \lr{NOT} (که با خط بالای متغیر نشان می‌دهیم: $\overline{x}$) ترکیب می‌شوند. یک \textbf{فرمول بولی} عبارتی از متغیرهای بولی و این عملگرهاست، مثلاً
\[\phi=(\overline{x}\wedge y)\vee(x\wedge\overline{z}).\]
فرمول بولی \textbf{صدق‌پذیر} است اگر مقداردهی‌ای از $0$ و $1$ به متغیرها وجود داشته باشد که مقدار فرمول را $1$ کند. مثلاً $\phi$ با $x=0$، $y=1$ و $z=0$ صدق‌پذیر است.
\[\prob{SAT}=\{\enc{\phi}\st\phi\ \text{\rl{یک فرمول بولی صدق‌پذیر است}}\}.\]

\begin{theorem}[کوک–لوین]
$\prob{SAT}$ NP-تمام است.
\end{theorem}

\begin{proofidea}
\begin{itemize}
\item $\prob{SAT}\in\cNP$ ساده است: گواه یک مقداردهی صدق‌کننده است که در زمان چندجمله‌ای بررسی می‌شود.
\item بخش اصلی: باید نشان دهیم \textbf{هر} زبان $A\in\cNP$ به‌طور چندجمله‌ای به $\prob{SAT}$ کاهش می‌یابد. فرض کنید $N$ یک ماشین غیرقطعی با زمان چندجمله‌ای برای $A$ باشد. برای هر رشته $w$ یک فرمول بولی $\phi$ می‌سازیم که \textbf{محاسبه $N$ روی $w$ را شبیه‌سازی کند}: اگر $N$ بپذیرد، $\phi$ صدق‌پذیر است و اگر $N$ رد کند، نیست.
\end{itemize}
\end{proofidea}

\begin{proof}
\textbf{$\prob{SAT}\in\cNP$:} ماشین غیرقطعی یک مقداردهی به متغیرها حدس می‌زند و فرمول را با آن ارزیابی می‌کند؛ همه این کارها در زمان چندجمله‌ای انجام می‌شود.

\textbf{ساخت کاهش.} فرض کنید $N$ یک ماشین غیرقطعی باشد که $A$ را در زمان $n^k$ تصمیم می‌گیرد (برای سادگی زمان را $n^k$ می‌گیریم؛ با یک ثابت کار تغییر نمی‌کند). یک \textbf{جدول} (\lr{tableau}) برای $N$ روی $w$، جدولی $n^k\times n^k$ است که سطرهایش پیکربندی‌های یک شاخه از محاسبه $N$ روی $w$ هستند: سطر اول پیکربندی شروع و هر سطر پیکربندی بعدی سطر قبل است. برای راحتی هر سطر را با نماد $\#$ شروع و تمام می‌کنیم.
\begin{ltrfig}
\begin{tikzpicture}[font=\scriptsize,scale=0.9]
  \draw[thick,jzMain] (0,0) rectangle (6,5);
  \foreach \y in {4.5,4.0,3.5} \draw[jzMain!50] (0,\y) -- (6,\y);
  \draw[jzMain!50] (0,0.5) -- (6,0.5);
  \foreach \x in {0.5,5.5} \draw[jzMain!50] (\x,0) -- (\x,5);
  \foreach \x/\s in {0.25/\#,0.75/q_0,1.25/w_1,1.75/w_2,2.5/\cdots,3.25/w_n,3.75/\sqcup,4.6/\cdots,5.75/\#} \node at (\x,4.75) {$\s$};
  \foreach \x in {1.0,1.5,2.0,3.0,3.5,4.0} \draw[jzMain!50] (\x,4.5) -- (\x,5);
  \node at (0.25,4.25) {$\#$}; \node at (5.75,4.25) {$\#$}; \node at (0.25,0.25) {$\#$}; \node at (5.75,0.25) {$\#$};
  \draw[jzAccent,thick] (2.5,2.0) rectangle (4.0,3.0); \draw[jzAccent] (2.5,2.5) -- (4.0,2.5); \draw[jzAccent] (3.0,2.0) -- (3.0,3.0); \draw[jzAccent] (3.5,2.0) -- (3.5,3.0);
  \node[anchor=west] at (6.2,4.75) {start configuration};
  \node[anchor=west] at (6.2,4.25) {second configuration};
  \node[anchor=west] at (6.2,0.25) {$n^k$-th configuration};
  \node[anchor=west] (wl) at (6.2,2.5) {window}; \draw[->,jzAccent] (wl.west) -- (4.0,2.5);
  \draw[<->] (-0.4,0) -- node[left] {$n^k$} (-0.4,5);
  \draw[<->] (0,-0.4) -- node[below] {$n^k$} (6,-0.4);
\end{tikzpicture}
\end{ltrfig}
هر جدول پذیرنده متناظر با یک شاخه پذیرنده از محاسبه $N$ روی $w$ است (جدولی که یکی از سطرهایش پیکربندی پذیرش باشد). پس مسئله «آیا $N$ رشته $w$ را می‌پذیرد؟» معادل این است که «آیا یک جدول پذیرنده برای $N$ روی $w$ وجود دارد؟».

\textbf{متغیرها.} فرض کنید $Q$ و $\Gamma$ مجموعه حالت‌ها و الفبای نوار $N$ باشند و $C=Q\cup\Gamma\cup\{\#\}$. برای هر خانه $(i,j)$ جدول و هر نماد $s\in C$ یک متغیر بولی $x_{i,j,s}$ تعریف می‌کنیم با این معنی:
\[x_{i,j,s}=1\iff\text{\rl{خانه $(i,j)$ جدول نماد $s$ را دارد}}.\]
فرمول $\phi$ از چهار بخش تشکیل می‌شود:
\[\phi=\phi_{\mathrm{cell}}\wedge\phi_{\mathrm{start}}\wedge\phi_{\mathrm{move}}\wedge\phi_{\mathrm{accept}}.\]

\textbf{۱. هر خانه دقیقاً یک نماد دارد:}
\[\phi_{\mathrm{cell}}=\bigwedge_{1\le i,j\le n^k}\left[\Big(\bigvee_{s\in C}x_{i,j,s}\Big)\wedge\Big(\bigwedge_{\substack{s,t\in C\\ s\neq t}}(\overline{x_{i,j,s}}\vee\overline{x_{i,j,t}})\Big)\right].\]
بخش اول داخل کروشه می‌گوید «دست‌کم یک نماد در خانه هست» و بخش دوم «بیش از یک نماد نیست». این شرط برای هر خانه جدول برقرار است.

\textbf{۲. سطر اول پیکربندی شروع است:}
\[\phi_{\mathrm{start}}=x_{1,1,\#}\wedge x_{1,2,q_0}\wedge x_{1,3,w_1}\wedge x_{1,4,w_2}\wedge\cdots\wedge x_{1,n+2,w_n}\wedge x_{1,n+3,\sqcup}\wedge\cdots\wedge x_{1,n^k-1,\sqcup}\wedge x_{1,n^k,\#}.\]

\textbf{۳. جدول پیکربندی پذیرش دارد:}
\[\phi_{\mathrm{accept}}=\bigvee_{1\le i,j\le n^k}x_{i,j,\qacc}.\]

\textbf{۴. هر سطر به‌طور مجاز از سطر قبل نتیجه می‌شود.} این را با بررسی \textbf{پنجره‌های} $2\times3$ از خانه‌ها انجام می‌دهیم. یک پنجره \textbf{مجاز} است اگر با تابع گذار $N$ تناقض نداشته باشد. مثلاً اگر $\delta(q_1,a)=\{(q_1,b,R)\}$ و $\delta(q_1,b)=\{(q_2,c,L),(q_2,a,R)\}$، پنجره‌های زیر مجازند:
\begin{ltrfig}
\newcommand{\win}[6]{\begin{tabular}{|c|c|c|}\hline $#1$ & $#2$ & $#3$\\ \hline $#4$ & $#5$ & $#6$\\ \hline\end{tabular}}
\begin{tabular}{cccccc}
(a)\ \win{a}{q_1}{b}{q_2}{a}{c} & \quad (b)\ \win{a}{q_1}{b}{a}{a}{q_2} & \quad (c)\ \win{a}{a}{q_1}{a}{a}{b}\\[10pt]
(d)\ \win{\#}{b}{a}{\#}{b}{a} & \quad (e)\ \win{a}{b}{a}{a}{b}{q_2} & \quad (f)\ \win{b}{b}{b}{c}{b}{b}
\end{tabular}
\end{ltrfig}
(در (a) و (b) گذار $\delta(q_1,b)$ به دو روش مختلف اجرا شده است. در (c) حالت $q_1$ در لبه راست پنجره است و نمادی که می‌خواند بیرون پنجره است؛ اگر آن نماد $a$ باشد، با $b$ عوض می‌شود و هد به راست می‌رود، پس جای $q_1$ نماد $b$ می‌نشیند؛ در (e) هد از سمت راست وارد پنجره شده؛ در (f) هد کنار پنجره است و نماد زیر آن عوض شده.) و پنجره‌های زیر مجاز نیستند:
\begin{ltrfig}
\newcommand{\win}[6]{\begin{tabular}{|c|c|c|}\hline $#1$ & $#2$ & $#3$\\ \hline $#4$ & $#5$ & $#6$\\ \hline\end{tabular}}
\begin{tabular}{ccc}
(a)\ \win{a}{b}{a}{a}{a}{a} & \quad (b)\ \win{a}{q_1}{b}{q_2}{a}{a} & \quad (c)\ \win{b}{q_1}{b}{q_2}{b}{q_2}
\end{tabular}
\end{ltrfig}
(در (a) نماد وسط بدون حضور هد عوض شده؛ در (b) گذار با تابع گذار نمی‌خواند؛ در (c) دو حالت در یک سطر آمده است.)

می‌توان ثابت کرد که اگر سطر اول پیکربندی شروع باشد و \textbf{همه} پنجره‌های جدول مجاز باشند، هر سطر پیکربندی‌ای است که به‌طور مجاز از سطر قبلی‌اش نتیجه می‌شود. (دلیل: هر خانه‌ای در سطر بالا که کنار هد نیست، در پنجره‌ای قرار دارد که آن را به‌طور تغییرنیافته به سطر پایین منتقل می‌کند، و پنجره‌هایی که هد را دارند، دقیقاً مطابق یک گذار $N$ به‌روز می‌شوند.) پنجره $(i,j)$ را پنجره‌ای می‌گیریم که خانه $(i,j)$ وسط سطر بالای آن است:
\[\phi_{\mathrm{move}}=\bigwedge_{1\le i<n^k,\ 1<j<n^k}\big(\text{\rl{پنجره $(i,j)$ مجاز است}}\big),\]
که «پنجره $(i,j)$ مجاز است» با فرمول زیر بیان می‌شود:
\[\bigvee_{\substack{a_1,\dots,a_6\\ \text{\rl{پنجره مجاز}}}}\big(x_{i,j-1,a_1}\wedge x_{i,j,a_2}\wedge x_{i,j+1,a_3}\wedge x_{i+1,j-1,a_4}\wedge x_{i+1,j,a_5}\wedge x_{i+1,j+1,a_6}\big).\]

\textbf{درستی.} اگر $N$ رشته $w$ را بپذیرد، جدول پذیرنده‌ای هست و مقداردهی متناظر با آن همه چهار بخش را صادق می‌کند. برعکس، هر مقداردهی صدق‌کننده $\phi$ یک جدول را مشخص می‌کند (طبق $\phi_{\mathrm{cell}}$)، که با پیکربندی شروع آغاز می‌شود ($\phi_{\mathrm{start}}$)، هر سطرش به‌طور مجاز از قبلی نتیجه می‌شود ($\phi_{\mathrm{move}}$) و به پذیرش می‌رسد ($\phi_{\mathrm{accept}}$). پس $N$ رشته $w$ را می‌پذیرد.

\textbf{زمان کاهش.} اندازه $\phi$ را تخمین می‌زنیم.
\begin{itemize}
\item جدول $n^k\times n^k=n^{2k}$ خانه دارد و هر خانه $l=\abs{C}$ متغیر دارد. $l$ به $N$ بستگی دارد و نه به ورودی، پس ثابت است. تعداد کل متغیرها $O(n^{2k})$ است.
\item $\phi_{\mathrm{cell}}$: برای هر خانه بخشی با اندازه ثابت؛ کل $O(n^{2k})$.
\item $\phi_{\mathrm{start}}$: برای هر خانه سطر اول بخشی ثابت؛ کل $O(n^k)$.
\item $\phi_{\mathrm{accept}}$ و $\phi_{\mathrm{move}}$: برای هر خانه بخشی ثابت (تعداد پنجره‌های مجاز ثابت است)؛ کل $O(n^{2k})$.
\item هر متغیر یک اندیس دارد که با $O(\log n)$ نماد نوشته می‌شود؛ پس طول کل $\phi$ برابر $O(n^{2k}\log n)$ است.
\end{itemize}
$\phi$ ساختار بسیار منظمی دارد و می‌توان آن را به‌سادگی در زمانی متناسب با طولش تولید کرد. پس کاهش چندجمله‌ای است و $A\le_P\prob{SAT}$.
\end{proof}

\section{3SAT}

\begin{definition}[فرمول \lr{cnf}]
\begin{itemize}
\item یک \textbf{لیترال} یک متغیر بولی یا نقیض آن است، مثل $x$ یا $\overline{x}$.
\item یک \textbf{بند} (\lr{clause}) چند لیترال است که با $\vee$ به هم وصل شده‌اند، مثل $(x_1\vee\overline{x_2}\vee\overline{x_3}\vee x_4)$.
\item یک فرمول بولی در \textbf{فرم نرمال عطفی} (\lr{cnf}) است اگر از چند بند تشکیل شده باشد که با $\wedge$ به هم وصل شده‌اند:
\[(x_1\vee\overline{x_2}\vee\overline{x_3}\vee x_4)\wedge(x_3\vee\overline{x_5}\vee x_6)\wedge(x_3\vee\overline{x_6}).\]
\item یک \textbf{فرمول 3cnf} فرمولی \lr{cnf} است که هر بندش دقیقاً سه لیترال دارد:
\[(x_1\vee\overline{x_2}\vee\overline{x_3})\wedge(x_3\vee\overline{x_5}\vee x_6)\wedge(x_3\vee\overline{x_6}\vee x_4)\wedge(x_4\vee x_5\vee x_6).\]
\end{itemize}
\end{definition}

\[\prob{3SAT}=\{\enc{\phi}\st\phi\ \text{\rl{یک فرمول 3cnf صدق‌پذیر است}}\}.\]

\begin{theorem}
$\prob{3SAT}$ NP-تمام است.
\end{theorem}
\begin{proof}
$\prob{3SAT}\in\cNP$ دقیقاً مثل $\prob{SAT}$ است. برای بخش دوم، کاهش کوک–لوین را طوری تغییر می‌دهیم که مستقیماً یک فرمول 3cnf بسازد.

\textbf{گام ۱: رساندن به فرم \lr{cnf}.} $\phi_{\mathrm{cell}}$، $\phi_{\mathrm{start}}$ و $\phi_{\mathrm{accept}}$ از قبل \lr{cnf} هستند ($\phi_{\mathrm{accept}}$ یک بند بزرگ است و $\phi_{\mathrm{start}}$ عطف بندهای تک‌لیترالی). $\phi_{\mathrm{move}}$ عطفی از فرمول‌های «پنجره $(i,j)$ مجاز است» است که هر کدام \textbf{فصل عطف‌ها} با اندازه ثابت است. هر فرمول با اندازه ثابت را می‌توان با قوانین پخش‌پذیری به \lr{cnf} با اندازه ثابت (شاید بزرگ‌تر، اما همچنان ثابت و مستقل از $n$) تبدیل کرد. پس کل فرمول \lr{cnf} می‌شود و اندازه‌اش فقط در یک ضریب ثابت بزرگ می‌شود.

\textbf{گام ۲: رساندن هر بند به دقیقاً سه لیترال.}
\begin{itemize}
\item \textbf{بندهای کوتاه:} یک لیترال را تکرار می‌کنیم: $(a_1\vee a_2)\ \to\ (a_1\vee a_2\vee a_2)$ و $(a_1)\to(a_1\vee a_1\vee a_1)$. این کار معنی بند را تغییر نمی‌دهد.
\item \textbf{بندهای بلند:} بند $(a_1\vee a_2\vee\cdots\vee a_l)$ با $l>3$ را با متغیرهای کمکی جدید $z_1,\dots,z_{l-3}$ به زنجیره زیر تبدیل می‌کنیم:
\[(a_1\vee a_2\vee z_1)\wedge(\overline{z_1}\vee a_3\vee z_2)\wedge(\overline{z_2}\vee a_4\vee z_3)\wedge\cdots\wedge(\overline{z_{l-3}}\vee a_{l-1}\vee a_l).\]
\end{itemize}
چرا تبدیل بندهای بلند صدق‌پذیری را حفظ می‌کند؟ اگر یک مقداردهی بند اصلی را صادق کند، یک لیترال $a_r$ مقدار $1$ دارد. قرار دهید $z_1=\cdots=z_{r-2}=1$ و بقیه $z$ها $=0$. آنگاه بندهای قبل از بند شامل $a_r$ با $z$ مثبتشان، بند شامل $a_r$ با خود $a_r$ و بندهای بعد از آن با $\overline{z}$ شان صادق می‌شوند. برعکس، اگر همه $a_i$ها $0$ باشند، زنجیره نمی‌تواند صادق باشد: بند اول $z_1=1$ را لازم دارد، بند دوم سپس $z_2=1$ و \dots\ تا بند آخر که با $z_{l-3}=1$ و $a_{l-1}=a_l=0$ کاذب می‌شود. پس زنجیره صدق‌پذیر است اگر و تنها اگر بند اصلی صدق‌پذیر باشد.

هر دو گام در زمان چندجمله‌ای انجام می‌شوند و اندازه فرمول را فقط چندجمله‌ای زیاد می‌کنند. پس هر زبان $\cNP$ به‌طور چندجمله‌ای به $\prob{3SAT}$ کاهش می‌یابد و $\prob{3SAT}$ NP-تمام است.
\end{proof}

\begin{note}
با داشتن $\prob{3SAT}$ و قضیه \ref{thm:npc-chain}، NP-تمام بودن مسئله‌های بسیار زیادی (مثل \lr{CLIQUE}، \lr{VERTEX-COVER}، \lr{HAMPATH} و \lr{SUBSET-SUM}) با کاهش از $\prob{3SAT}$ ثابت می‌شود. امروز هزاران مسئله NP-تمام می‌شناسیم.
\end{note}

\begin{summary}
\begin{itemize}
\item $A\le_PB$: کاهش نگاشتی با تابعی که در زمان چندجمله‌ای محاسبه می‌شود. اگر $B\in\cP$ آنگاه $A\in\cP$.
\item NP-تمام: در $\cNP$ است و همه $\cNP$ به آن کاهش می‌یابد. اگر یکی در $\cP$ باشد، $\cP=\cNP$.
\item کوک–لوین: $\prob{SAT}$ NP-تمام است؛ جدول $n^k\times n^k$ و فرمول $\phi_{\mathrm{cell}}\wedge\phi_{\mathrm{start}}\wedge\phi_{\mathrm{move}}\wedge\phi_{\mathrm{accept}}$ با پنجره‌های $2\times3$.
\item $\prob{3SAT}$ هم NP-تمام است؛ با تبدیل به \lr{cnf} و شکستن بندهای بلند با متغیرهای کمکی.
\end{itemize}
\end{summary}

\vspace{1cm}
\begin{center}
\textit{«از آغاز شروع کن و ادامه بده تا به پایان برسی؛ آنجا بایست.»}\\[2pt]
{\small — لوئیس کارول، آلیس در سرزمین عجایب}
\end{center}
